% Options for packages loaded elsewhere
% Options for packages loaded elsewhere
\PassOptionsToPackage{unicode}{hyperref}
\PassOptionsToPackage{hyphens}{url}
\PassOptionsToPackage{dvipsnames,svgnames,x11names}{xcolor}
%
\documentclass[
  letterpaper,
  DIV=11,
  numbers=noendperiod]{scrartcl}
\usepackage{xcolor}
\usepackage{amsmath,amssymb}
\setcounter{secnumdepth}{-\maxdimen} % remove section numbering
\usepackage{iftex}
\ifPDFTeX
  \usepackage[T1]{fontenc}
  \usepackage[utf8]{inputenc}
  \usepackage{textcomp} % provide euro and other symbols
\else % if luatex or xetex
  \usepackage{unicode-math} % this also loads fontspec
  \defaultfontfeatures{Scale=MatchLowercase}
  \defaultfontfeatures[\rmfamily]{Ligatures=TeX,Scale=1}
\fi
\usepackage{lmodern}
\ifPDFTeX\else
  % xetex/luatex font selection
\fi
% Use upquote if available, for straight quotes in verbatim environments
\IfFileExists{upquote.sty}{\usepackage{upquote}}{}
\IfFileExists{microtype.sty}{% use microtype if available
  \usepackage[]{microtype}
  \UseMicrotypeSet[protrusion]{basicmath} % disable protrusion for tt fonts
}{}
\makeatletter
\@ifundefined{KOMAClassName}{% if non-KOMA class
  \IfFileExists{parskip.sty}{%
    \usepackage{parskip}
  }{% else
    \setlength{\parindent}{0pt}
    \setlength{\parskip}{6pt plus 2pt minus 1pt}}
}{% if KOMA class
  \KOMAoptions{parskip=half}}
\makeatother
% Make \paragraph and \subparagraph free-standing
\makeatletter
\ifx\paragraph\undefined\else
  \let\oldparagraph\paragraph
  \renewcommand{\paragraph}{
    \@ifstar
      \xxxParagraphStar
      \xxxParagraphNoStar
  }
  \newcommand{\xxxParagraphStar}[1]{\oldparagraph*{#1}\mbox{}}
  \newcommand{\xxxParagraphNoStar}[1]{\oldparagraph{#1}\mbox{}}
\fi
\ifx\subparagraph\undefined\else
  \let\oldsubparagraph\subparagraph
  \renewcommand{\subparagraph}{
    \@ifstar
      \xxxSubParagraphStar
      \xxxSubParagraphNoStar
  }
  \newcommand{\xxxSubParagraphStar}[1]{\oldsubparagraph*{#1}\mbox{}}
  \newcommand{\xxxSubParagraphNoStar}[1]{\oldsubparagraph{#1}\mbox{}}
\fi
\makeatother

\usepackage{color}
\usepackage{fancyvrb}
\newcommand{\VerbBar}{|}
\newcommand{\VERB}{\Verb[commandchars=\\\{\}]}
\DefineVerbatimEnvironment{Highlighting}{Verbatim}{commandchars=\\\{\}}
% Add ',fontsize=\small' for more characters per line
\usepackage{framed}
\definecolor{shadecolor}{RGB}{241,243,245}
\newenvironment{Shaded}{\begin{snugshade}}{\end{snugshade}}
\newcommand{\AlertTok}[1]{\textcolor[rgb]{0.68,0.00,0.00}{#1}}
\newcommand{\AnnotationTok}[1]{\textcolor[rgb]{0.37,0.37,0.37}{#1}}
\newcommand{\AttributeTok}[1]{\textcolor[rgb]{0.40,0.45,0.13}{#1}}
\newcommand{\BaseNTok}[1]{\textcolor[rgb]{0.68,0.00,0.00}{#1}}
\newcommand{\BuiltInTok}[1]{\textcolor[rgb]{0.00,0.23,0.31}{#1}}
\newcommand{\CharTok}[1]{\textcolor[rgb]{0.13,0.47,0.30}{#1}}
\newcommand{\CommentTok}[1]{\textcolor[rgb]{0.37,0.37,0.37}{#1}}
\newcommand{\CommentVarTok}[1]{\textcolor[rgb]{0.37,0.37,0.37}{\textit{#1}}}
\newcommand{\ConstantTok}[1]{\textcolor[rgb]{0.56,0.35,0.01}{#1}}
\newcommand{\ControlFlowTok}[1]{\textcolor[rgb]{0.00,0.23,0.31}{\textbf{#1}}}
\newcommand{\DataTypeTok}[1]{\textcolor[rgb]{0.68,0.00,0.00}{#1}}
\newcommand{\DecValTok}[1]{\textcolor[rgb]{0.68,0.00,0.00}{#1}}
\newcommand{\DocumentationTok}[1]{\textcolor[rgb]{0.37,0.37,0.37}{\textit{#1}}}
\newcommand{\ErrorTok}[1]{\textcolor[rgb]{0.68,0.00,0.00}{#1}}
\newcommand{\ExtensionTok}[1]{\textcolor[rgb]{0.00,0.23,0.31}{#1}}
\newcommand{\FloatTok}[1]{\textcolor[rgb]{0.68,0.00,0.00}{#1}}
\newcommand{\FunctionTok}[1]{\textcolor[rgb]{0.28,0.35,0.67}{#1}}
\newcommand{\ImportTok}[1]{\textcolor[rgb]{0.00,0.46,0.62}{#1}}
\newcommand{\InformationTok}[1]{\textcolor[rgb]{0.37,0.37,0.37}{#1}}
\newcommand{\KeywordTok}[1]{\textcolor[rgb]{0.00,0.23,0.31}{\textbf{#1}}}
\newcommand{\NormalTok}[1]{\textcolor[rgb]{0.00,0.23,0.31}{#1}}
\newcommand{\OperatorTok}[1]{\textcolor[rgb]{0.37,0.37,0.37}{#1}}
\newcommand{\OtherTok}[1]{\textcolor[rgb]{0.00,0.23,0.31}{#1}}
\newcommand{\PreprocessorTok}[1]{\textcolor[rgb]{0.68,0.00,0.00}{#1}}
\newcommand{\RegionMarkerTok}[1]{\textcolor[rgb]{0.00,0.23,0.31}{#1}}
\newcommand{\SpecialCharTok}[1]{\textcolor[rgb]{0.37,0.37,0.37}{#1}}
\newcommand{\SpecialStringTok}[1]{\textcolor[rgb]{0.13,0.47,0.30}{#1}}
\newcommand{\StringTok}[1]{\textcolor[rgb]{0.13,0.47,0.30}{#1}}
\newcommand{\VariableTok}[1]{\textcolor[rgb]{0.07,0.07,0.07}{#1}}
\newcommand{\VerbatimStringTok}[1]{\textcolor[rgb]{0.13,0.47,0.30}{#1}}
\newcommand{\WarningTok}[1]{\textcolor[rgb]{0.37,0.37,0.37}{\textit{#1}}}

\usepackage{longtable,booktabs,array}
\newcounter{none} % for unnumbered tables
\usepackage{calc} % for calculating minipage widths
% Correct order of tables after \paragraph or \subparagraph
\usepackage{etoolbox}
\makeatletter
\patchcmd\longtable{\par}{\if@noskipsec\mbox{}\fi\par}{}{}
\makeatother
% Allow footnotes in longtable head/foot
\IfFileExists{footnotehyper.sty}{\usepackage{footnotehyper}}{\usepackage{footnote}}
\makesavenoteenv{longtable}
\usepackage{graphicx}
\makeatletter
\newsavebox\pandoc@box
\newcommand*\pandocbounded[1]{% scales image to fit in text height/width
  \sbox\pandoc@box{#1}%
  \Gscale@div\@tempa{\textheight}{\dimexpr\ht\pandoc@box+\dp\pandoc@box\relax}%
  \Gscale@div\@tempb{\linewidth}{\wd\pandoc@box}%
  \ifdim\@tempb\p@<\@tempa\p@\let\@tempa\@tempb\fi% select the smaller of both
  \ifdim\@tempa\p@<\p@\scalebox{\@tempa}{\usebox\pandoc@box}%
  \else\usebox{\pandoc@box}%
  \fi%
}
% Set default figure placement to htbp
\def\fps@figure{htbp}
\makeatother





\setlength{\emergencystretch}{3em} % prevent overfull lines

\providecommand{\tightlist}{%
  \setlength{\itemsep}{0pt}\setlength{\parskip}{0pt}}



 


% Page furniture for the project PDFs, ported from the exam-documentclass originals the 2024 bootcamp used before the Quarto migration.
%
% Pulled in by every project's YAML:
%
%   format:
%     pdf:
%       include-in-header: header.tex
%     latex:
%       include-in-header: header.tex
%
% Both keys are needed: pdf and latex are separate formats to Quarto, and with the key on pdf alone the shipped .tex comes out without the header its own .pdf has.
%
% The exam class provided \header, \headrule and \footer; Quarto's pdf output uses the default article class, so fancyhdr supplies the same furniture.
% The exam class's \numpages becomes \pageref*{LastPage}, which needs two passes --- Quarto runs latexmk, so that is already handled.
% The star matters: plain \pageref is a hyperref link, so the page count comes out link-blue.
%
% Editing this file changes every project PDF, so it is folded into deploy.sh's GLOBAL_HASH: a change here rebuilds the whole site rather than serving stale PDFs.
%
% ---------------------------------------------------------------------------
% CURRENTLY INERT: every rule below is commented out, matching the P230A repo
% this was ported from. Uncomment the block to turn the running header on.
% ADAPTING THIS FOR ANOTHER COURSE: change the two strings below, nothing else.
% ---------------------------------------------------------------------------

% \usepackage{fancyhdr}
% \usepackage{lastpage}

% \pagestyle{fancy}
% \fancyhf{}

% \fancyhead[L]{MCSB Bootcamp: Mathematical and Computational Track}  % <- course name
% \fancyfoot[L]{\textit{jun.allard@uci.edu}}                          % <- contact

% \fancyfoot[R]{Page \thepage\ of \pageref*{LastPage}}

% \renewcommand{\headrulewidth}{0.4pt}
% \renewcommand{\footrulewidth}{0pt}

% % Quarto's title block starts the document with \maketitle, which forces the plain pagestyle on the first page and would otherwise drop the header exactly where it is most wanted.
% \fancypagestyle{plain}{%
%   \fancyhf{}%
%   \fancyhead[L]{MCSB Bootcamp: Mathematical and Computational Track}%
%   \fancyfoot[L]{\textit{jun.allard@uci.edu}}%
%   \fancyfoot[R]{Page \thepage\ of \pageref*{LastPage}}%
%   \renewcommand{\headrulewidth}{0.4pt}%
%   \renewcommand{\footrulewidth}{0pt}%
% }
\KOMAoption{captions}{tableheading}
\makeatletter
\@ifpackageloaded{caption}{}{\usepackage{caption}}
\AtBeginDocument{%
\ifdefined\contentsname
  \renewcommand*\contentsname{Table of contents}
\else
  \newcommand\contentsname{Table of contents}
\fi
\ifdefined\listfigurename
  \renewcommand*\listfigurename{List of Figures}
\else
  \newcommand\listfigurename{List of Figures}
\fi
\ifdefined\listtablename
  \renewcommand*\listtablename{List of Tables}
\else
  \newcommand\listtablename{List of Tables}
\fi
\ifdefined\figurename
  \renewcommand*\figurename{Figure}
\else
  \newcommand\figurename{Figure}
\fi
\ifdefined\tablename
  \renewcommand*\tablename{Table}
\else
  \newcommand\tablename{Table}
\fi
}
\@ifpackageloaded{float}{}{\usepackage{float}}
\floatstyle{ruled}
\@ifundefined{c@chapter}{\newfloat{codelisting}{h}{lop}}{\newfloat{codelisting}{h}{lop}[chapter]}
\floatname{codelisting}{Listing}
\newcommand*\listoflistings{\listof{codelisting}{List of Listings}}
\makeatother
\makeatletter
\makeatother
\makeatletter
\@ifpackageloaded{caption}{}{\usepackage{caption}}
\@ifpackageloaded{subcaption}{}{\usepackage{subcaption}}
\makeatother
\usepackage{bookmark}
\IfFileExists{xurl.sty}{\usepackage{xurl}}{} % add URL line breaks if available
\urlstyle{same}
\hypersetup{
  pdftitle={Project 5: FitzHugh-Nagumo and excitability},
  pdfauthor={Jun Allard},
  colorlinks=true,
  linkcolor={blue},
  filecolor={Maroon},
  citecolor={Blue},
  urlcolor={Blue},
  pdfcreator={LaTeX via pandoc}}

\ifXeTeX
\special{pdf:trailerid [ <64c529d4ea1052dffcbe2f12e828a3bf> <64c529d4ea1052dffcbe2f12e828a3bf> ]}
\fi
\ifPDFTeX
\pdftrailerid{}
\pdftrailer{/ID [ <64c529d4ea1052dffcbe2f12e828a3bf> <64c529d4ea1052dffcbe2f12e828a3bf> ]}
\fi
\ifLuaTeX
\pdfvariable trailerid {[ <64c529d4ea1052dffcbe2f12e828a3bf> <64c529d4ea1052dffcbe2f12e828a3bf> ]}
\fi


\title{Project 5: FitzHugh-Nagumo and excitability}
\usepackage{etoolbox}
\makeatletter
\providecommand{\subtitle}[1]{% add subtitle to \maketitle
  \apptocmd{\@title}{\par {\large #1 \par}}{}{}
}
\makeatother
\subtitle{MCSB Bootcamp --- Mathematical and Computational Track}
\author{Jun Allard}
\date{}
\begin{document}
\maketitle


\subsection{Background, model, and Part
I}\label{background-model-and-part-i}

The phenomenon of excitability exists in many biological systems,
including in the electrophysiology of neurons. The FitzHugh-Nagumo
equations

\[
\begin{aligned}
\frac{dv}{dt} &= v - \frac{1}{3} v^3 - w \\
\frac{dw}{dt} &= \epsilon \left( v + a - bw\right)
\end{aligned}
\]

describe neuron electrophysiology where, roughly speaking, \(v\) is the
electrical potential (voltage) across the cell's membrane, and \(w\) is
the activity of ion pumps. The parameters \(\epsilon\), \(a\) and \(b\)
represent properties of the ion pumps. The model has been
nondimensionalized. Both \(v\) and \(w\) can be negative or positive.

\begin{enumerate}
\def\labelenumi{\alph{enumi}.}
\tightlist
\item
  Confirm that for \(\epsilon=0.08\), \(a=0.5\), \(b=0.2\), the system
  is oscillatory.
\item
  Confirm that for \(\epsilon=0.08\), \(a=1.0\), \(b=0.2\), the system
  has the following property: if you choose initial conditions
  \(v(0)=-1.5\), \(w(0)=-0.5\), the system evolves directly towards a
  stable steady state, but if you choose initial conditions
  \(v(0)=0.0\), \(w(0)=-0.5\), the system moves away from the steady
  state before eventually converging towards it. This behavior is called
  \emph{excitability}. At these parameters, what is the steady state
  value of \(v\) and \(w\)?
\end{enumerate}

\subsection{Part II: Current injection}\label{part-ii-current-injection}

Assume the neuron is at rest (at its steady state), and another cell
injects a current into it. The current is injected between \(t=40\) and
\(t=47\), and has a strength of \(I_0=1.0\). In the model, this means

\[
\begin{aligned}
\frac{dv}{dt} &= v - \frac{1}{3} v^3 - w + I(t)\\
\frac{dw}{dt} &= \epsilon \left( v + a - bw\right)
\end{aligned}
\]

where

\[
I(t) = \begin{cases}
I_0 & \qquad t_\mathrm{start} < t < t_\mathrm{stop}\\
0 & \qquad \text{otherwise}
\end{cases}
\]

or, in Matlab,

\begin{Shaded}
\begin{Highlighting}[]
\VariableTok{I0} \OperatorTok{=} \FloatTok{1.0}\OperatorTok{;}
\VariableTok{tStart} \OperatorTok{=} \FloatTok{40}\OperatorTok{;}
\VariableTok{tStop} \OperatorTok{=} \FloatTok{47}\OperatorTok{;}
\VariableTok{I} \OperatorTok{=} \OperatorTok{@}\NormalTok{(}\VariableTok{t}\NormalTok{) }\VariableTok{I0}\OperatorTok{*}\NormalTok{(}\VariableTok{t}\OperatorTok{\textgreater{}}\VariableTok{tStart}\NormalTok{)}\OperatorTok{.*}\NormalTok{(}\VariableTok{t}\OperatorTok{\textless{}}\VariableTok{tStop}\NormalTok{)}\OperatorTok{;}
\end{Highlighting}
\end{Shaded}

\begin{enumerate}
\def\labelenumi{\alph{enumi}.}
\setcounter{enumi}{2}
\tightlist
\item
  At the excitable parameters from above (\(a=1.0\)), simulate the
  system with initial conditions at the steady state (or very close),
  and simulate a 7-second injection starting at \(t=40\) as above.
\end{enumerate}

\subsection{Part III: Ring of ten
cells}\label{part-iii-ring-of-ten-cells}

Neurons are connected in a neural network. Suppose there are ten cells,
each with membrane potential and ion pump activity obeying the
FitzHugh-Nagumo equations for \(v_i(t)\) and \(w_i(t)\) where
\(i=1 \ldots 10\) indexes the cells. The cells are electrically
connected so that

\[
\begin{aligned}
\frac{dv_i}{dt} &= v_i - \frac{1}{3} v_i^3 - w_i + I_i(t) + D\left( v_{i-1} - 2 v_i + v_{i+1}\right)\\
\frac{dw_i}{dt} &= \epsilon \left( v_i + a - bw_i\right)
\end{aligned}
\]

where \(D=0.9\) is a new parameter that describes the electrical
connectivity of the neighboring cells. The ion pumps are not connected
between cells, so the \(w\) equation is unchanged. For simplicity, let's
assume the cells are connected in a ring, so that

\[\frac{dv_1}{dt} = v_1 - \frac{1}{3} v_1^3 - w_1 + I_1(t) + D\left( v_{10} - 2 v_1 + v_{2}\right)\]

and similarly for the tenth cell,

\[\frac{dv_{10}}{dt} = v_{10} - \frac{1}{3} v_{10}^3 - w_{10} + I_{10}(t) + D\left( v_{9} - 2 v_{10} + v_{1}\right).\]

\begin{enumerate}
\def\labelenumi{\alph{enumi}.}
\setcounter{enumi}{3}
\item
  Write code to simulate these ten cells. There will be 20 equations,
  \(v\) and \(w\) for each cell.

  \begin{enumerate}
  \def\labelenumii{\roman{enumii}.}
  \item
    Assume there is no injection current. We expect all ten cells to
    settle at the same steady state. Make two plots. First, plot a time
    series of the membrane potential of all ten cells as a function of
    time. Second, make a movie of the voltage in all ten cells, where
    the horizontal axis is the cell number.

\begin{Shaded}
\begin{Highlighting}[]
\CommentTok{\% movie}
\KeywordTok{for} \VariableTok{nt}\OperatorTok{=}\FloatTok{1}\OperatorTok{:}\VariableTok{numel}\NormalTok{(}\VariableTok{T}\NormalTok{)}
    \VariableTok{figure}\NormalTok{(}\FloatTok{5}\NormalTok{)}\OperatorTok{;} \VariableTok{clf}\OperatorTok{;} \VariableTok{hold} \VariableTok{on}\OperatorTok{;} \VariableTok{box} \VariableTok{on}\OperatorTok{;}
    \VariableTok{plot}\NormalTok{(}\VariableTok{X}\NormalTok{(}\VariableTok{nt}\OperatorTok{,}\FloatTok{1}\OperatorTok{:}\FloatTok{10}\NormalTok{))}\OperatorTok{;}
    \VariableTok{set}\NormalTok{(}\VariableTok{gca}\OperatorTok{,}\SpecialStringTok{\textquotesingle{}ylim\textquotesingle{}}\OperatorTok{,}\NormalTok{ [}\OperatorTok{{-}}\FloatTok{2.5}\OperatorTok{,}\FloatTok{2.5}\NormalTok{])}
    \VariableTok{xlabel}\NormalTok{(}\SpecialStringTok{\textquotesingle{}Cell\textquotesingle{}}\NormalTok{)}\OperatorTok{;}
    \VariableTok{ylabel}\NormalTok{(}\SpecialStringTok{\textquotesingle{}Voltage\textquotesingle{}}\NormalTok{)}
    \VariableTok{drawnow}\OperatorTok{;}
\KeywordTok{end}
\end{Highlighting}
\end{Shaded}
  \item
    Now assume that the fourth cell (and only the fourth cell) receives
    an injection current \(I(t)\) as above, between \(t=40\) and
    \(t=47\). Make a time series with all ten cells. Make a movie with
    the voltage for all cells.\footnote{This type of behavior is called
      an \emph{excitable traveling wave pulse}. These are unlike the
      harmonic pulses familiar from sound waves, vibrations, and light.
      For example, when two excitable pulses collide, they annihilate.}

    Hint: You do not need to make it generalize to different numbers of
    cells. In other words, don't try to make it work for an arbitrary
    number of cells and then set the number of cells to 10. Just make it
    work for 10 cells. (Unless you want to make it generalized.)
  \end{enumerate}
\end{enumerate}




\end{document}
